\section{Précédents dans les notes de Grothendieck}\label{sec:precedents}

Cette section rassemble deux précédents tirés des notes de Grothendieck. Les tabliers de ses notes figurent dans le registre (\S\ref{sec:grdecks}).

\subsection{Préordres et topos spatiaux: dossier 161-3, p.~26, et dossier 37, p.~48}

À la p.~26 d'une note manuscrite non datée sur les topos \cite{grothendieck-161-3}, Grothendieck énumère des conditions équivalentes portant sur un topos $E$: être équivalent à la fois au topos d'un espace et à un topos de préfaisceaux; être le topos d'un espace \q{essentiel} (un espace dont tous les points sont essentiels); être équivalent à $\widehat I$ pour un ensemble ordonné $I$; être engendré par les sous-objets de l'objet final et avoir suffisamment de points essentiels. Il demande de prouver que $\widehat C$ vérifie ces conditions si et seulement si $C$ est préordonnée, et que le topos d'un espace $X$ les vérifie si et seulement si les points essentiels de $X$ sont très denses; il note en outre des 2-équivalences entre topos spatiaux essentiels, préordres et espaces essentiels. Ce sont là des ponts au sens de cet article, lus sur le tablier des préordres avec pour arches les sous-objets de l'objet final et les points essentiels. Le générateur retrouve le premier exercice par l'arche de Morita: si $\widehat C\simeq\widehat P$ pour un préordre $P$, alors $\Kar(C)\simeq\Kar(P)\simeq P$, et $C$, sous-catégorie pleine de son enveloppe de Karoubi, est mince; réciproquement une catégorie mince est équivalente à son squelette mince, qui est un ordre (Lean: \lean{exists_preorder_moritaEquivalentCat_iff_isThin}). Le second exercice, sur les espaces, n'est pas formalisé ici. La même identification figure dans ses notes pour la réédition de SGA \cite[p.~48]{grothendieck-37}: pour un ensemble ordonné $I$, les ouverts de $\widehat I$ sont les cribles de $I$, d'où $\widehat I\simeq\mathrm{Top}(I)$.

\subsection{Sous-topos et quotients d'une théorie: dossier 161-2, p.~17 et 44}

Dans ses notes d'algèbre universelle \cite[p.~17]{grothendieck-161-2}, sous le titre \q{sous-théorie d'une $\Delta$-théorie (représentable par un topos) et sous-topos}, Grothendieck prend un topos $R$ et un ensemble $\Sigma$ de flèches de $R$: les morphismes qui inversent les flèches de $\Sigma$ se factorisent par un sous-topos $R'$ de $R$, et tout sous-topos de $R$ s'obtient ainsi. Lu avec pour $R$ le topos classifiant d'une théorie, c'est en substance la correspondance entre quotients d'une théorie et sous-topos de son topos classifiant sur laquelle Caramello fonde ses ponts \cite{caramello-tst}. À la p.~44 du même dossier, il lit des structures comme des ouverts d'un topos classifiant.
